\origpage{31}{GRAPH MINORS. XIV}
\noindent disjoint circuits, forming the boundaries of a nested sequence of discs all containing $u$, and that there is a large number of disjoint paths between the inner and outer circuits. In these circumstances a theorem of [3] can be applied to arrange that the end-vertices of the cross form a free set. There is only one other serious problem; it is natural to choose $\Lambda$ so that its boundary is a circuit of $H$, because then hypothesis (ii) provides information about $\Lambda$. But if we do so, then certain members of the 3-segregation of $G'$ may be arranged partly inside and partly outside $\Lambda$, and this causes confusion when we try to rearrange the part of the 3-segregation within $\Lambda$. To overcome this we use two discs, one bounded by a circuit of $H$, and the other, slightly larger, chosen so that no member of the 3-segregation is arranged partly inside and partly outside it.

The main part of the proof in this paper is concerned with proving the existence of the desired two discs. These lemmas are expressed in greater generality than is needed here, because the more general form (involving ``crooked transactions``) will be needed in the next paper of the series.

\sectionhead{5. Circuits in Cylinders}

We begin with some results about planar graphs. First, we need the following. A set $\mathscr{R}$ of regions of a drawing $H$ is \textit{connected} if for all $r, r' \in \mathscr{R}$ there is a sequence $r = r_1, r_2, \dots, r_k = r'$ of members of $\mathscr{R}$ such that for $1 \le i < k$, there is an edge of $H$ incident with $r_i$ and with $r_{i+1}$.

\stmt
(5.1) \textit{Let $H$ be a drawing in a surface $\Sigma$, and let $\Delta \subseteq \Sigma$ be a closed disc. Let $\mathscr{R}$ be a non-empty connected set of regions of $H$, such that $r \subseteq \Delta - \text{bd}(\Delta)$ for all $r \in \mathscr{R}$. Then there is a circuit $C$ of $H$ with $U(C) \subseteq \Delta$, bounding a disc in $\Delta$ including each member of $\mathscr{R}$, and such that every edge of $C$ is incident with a member of $\mathscr{R}$.}

This result is closely related to the elementary result that the minimal edge-cutsets of the dual graph of a planar graph correspond to the circuits of the graph, and may be proved similarly. \del{We omit further details.}

\fix{We give the details, using some results of [5].

\textit{Proof of (5.1).} We may assume that $\text{bd}(\Sigma) = \emptyset$. For otherwise, let $\Sigma'$ be the surface obtained from $\Sigma$ by pasting a closed disc onto each component of $\text{bd}(\Sigma)$. Then $H$ is a drawing in $\Sigma'$, and $\Delta$ is a closed disc in $\Sigma'$. Let $r \in \mathscr{R}$. Then $r$ is connected, and it is open in $\Sigma'$, since it is open in $\Sigma$ and disjoint from $\text{bd}(\Sigma)$ (because $r \subseteq \Delta - \text{bd}(\Delta)$). Moreover every point $x \in \bar{r} - r$ is in $U(H)$, for otherwise the region of $H$ in $\Sigma$ that contains $x$ is an open set meeting $r$, and so equals $r$, which is impossible since $x \notin r$. Hence $r$ is also a region of $H$ in $\Sigma'$. The incidences between the edges of $H$ and the members of $\mathscr{R}$ are the same in $\Sigma$ and in $\Sigma'$, and so (5.1) for $\Sigma'$ implies (5.1) for $\Sigma$.

Thus we assume that $\text{bd}(\Sigma) = \emptyset$. Let $N$ be the subdrawing of $H$ formed by the vertices and edges of $H$ that are incident both with a member of $\mathscr{R}$ and with a region of $H$ not in $\mathscr{R}$. (So $N = N_H(\mathscr{R})$ in the notation of [5].) Every vertex and every edge of $N$ is a subset of $\bar{r}$ for some $r \in \mathscr{R}$, and $\bar{r} \subseteq \Delta$ since $\Delta$ is closed. Hence $U(N) \subseteq \Delta$. For each circuit $C$ of $N$, $U(C)$ is an O-arc in the closed disc $\Delta$, and so it bounds a unique closed disc $\text{ins}(C) \subseteq \Delta$. Let $\theta = |E(N)|/2 + 1$. Every circuit of $N$ has length at most $|E(N)| < 2\theta$. We claim that $\text{ins}$ is a slope in $N$ of order $\theta$, in the sense of [5]. For the first slope axiom, let $C_1, C_2$ be circuits of $N$ with $U(C_1) \subseteq \text{ins}(C_2)$. Then $U(C_1)$ bounds a closed disc included in $\text{ins}(C_2)$, and this disc is $\text{ins}(C_1)$ by uniqueness, since $\text{ins}(C_2) \subseteq \Delta$. For the second, let $P_0, P_1, P_2$ be as in that axiom. If $U(P_0) \subseteq \text{ins}(P_1 \cup P_2)$, then $\text{ins}(P_1 \cup P_2)$ includes $U(P_2 \cup P_0)$ and $U(P_0 \cup P_1)$, and hence includes $\text{ins}(P_2 \cup P_0)$ and $\text{ins}(P_0 \cup P_1)$, by the first axiom. The same holds if $U(P_1) \subseteq \text{ins}(P_2 \cup P_0)$ or $U(P_2) \subseteq \text{ins}(P_0 \cup P_1)$, and so we may assume that none of these three inclusions holds. Since $U(P_0) \not\subseteq \text{ins}(P_1 \cup P_2)$ and $U(P_1) \not\subseteq \text{ins}(P_2 \cup P_0)$, the closed discs $\text{ins}(P_1 \cup P_2)$ and $\text{ins}(P_2 \cup P_0)$ meet only in $U(P_2)$, and their union is a closed disc bounded by $U(P_0 \cup P_1)$. This union is a subset of $\Delta$, and so it equals $\text{ins}(P_0 \cup P_1)$ by uniqueness. But it includes $U(P_2)$, a contradiction. This proves our claim. We apply the results of [5, Section 4] with $\Gamma = \Gamma' = N$. Since every circuit of $N$ has length $< 2\theta$, $N$ is confined.

\hspace*{12pt}(1) \textit{All the members of $\mathscr{R}$ are included in the same region of $N$.}

For let $r_1, r_2 \in \mathscr{R}$ be distinct, and both incident with an edge $e$ of $H$. Since $e$ is incident with at most two regions of $H$, it is incident with no region outside $\mathscr{R}$, and so $e \notin E(N)$. Choose $x \in e$. Then $x \notin U(N)$, since $e \notin E(N)$, and so $x$ belongs to a region $r'$ of $N$. Now $x \in \bar{r}_1 \cap \bar{r}_2$, and so $r'$ meets both $r_1$ and $r_2$. Since $r_1, r_2$ are connected and disjoint from $U(N)$, they are both included in $r'$. Since $\mathscr{R}$ is connected, (1) follows.

Let $r'$ be the region of $N$ that includes all the members of $\mathscr{R}$. Let $\text{ins}(N)$ be the union of $U(N)$ and of the discs $\text{ins}(C)$ for all circuits $C$ of $N$. By [5, (4.4)], $r_0 = \Sigma - \text{ins}(N)$ is a region of $N$. Since $\text{ins}(N) \subseteq \Delta$, it follows that $\Sigma - \Delta \subseteq r_0$. Now $\Sigma - \Delta$ is open, and it is non-empty, for otherwise $\Delta = \Sigma$, which is impossible since $\text{bd}(\Delta) \ne \emptyset$ and $\text{bd}(\Sigma) = \emptyset$. Since $U(H)$ includes no non-empty open subset of $\Sigma$, there is a point $y \in \Sigma - \Delta$ with $y \notin U(H)$. Let $r_2$ be the region of $H$ that contains $y$. Then $r_2 \notin \mathscr{R}$, since $y \notin \Delta$, and $r_2 \subseteq r_0$, since $r_2$ is connected, disjoint from $U(N)$, and contains $y \in r_0$.

\hspace*{12pt}(2) \textit{$r' \ne r_0$.}

For suppose that $r' = r_0$, and choose $r_1 \in \mathscr{R}$. Then $r_1, r_2 \subseteq r_0$, and so by [5, (5.1)], applied with $\Gamma = H$, $\Gamma' = N$ and $Z = \mathscr{R}$ (its proof does not use the slope), there are regions $r_3 \in \mathscr{R}$ and $r_4 \notin \mathscr{R}$ of $H$, both incident with some edge $e$ of $H$ with $e \notin E(N)$. But then $e \in E(N)$, from the definition of $N$, a contradiction. This proves (2).

Now $N$ does not capture $r_0$, since $y \in r_0$ and $y \notin \text{ins}(N)$. In the proof of [5, (4.4)] it is shown that a confined subdrawing captures at least one of any two distinct regions of it. By (2), it follows that $N$ captures $r'$, that is, there is a circuit $C$ of $N$ with $r' \subseteq \text{ins}(C)$. Then $U(C) \subseteq U(N) \subseteq \Delta$, and $U(C)$ bounds the disc $\text{ins}(C) \subseteq \Delta$, which includes each member of $\mathscr{R}$ by (1). Every edge of $C$ is an edge of $N$, and so is incident with a member of $\mathscr{R}$. This proves (5.1). $\square$}\fixnote{fix:5.1-proof-added}{The paper omits the proof of (5.1). Apart from the reduction in the first paragraph to a surface without boundary, the argument above is the one used in the formalization.}

\stmt
(5.2) \textit{Let $H$ be a drawing in $\Sigma$, and let $\Delta \subseteq \Sigma$ be a closed disc. Let $Z \subseteq \Delta - \text{bd}(\Delta)$ be non-empty and arc-wise connected, and such that $r \subseteq \Delta - \text{bd}(\Delta)$ for every region $r$ of $H$ with $\bar{r} \cap Z \ne \emptyset$. Then there is a circuit $C$ of $H$ with $U(C) \subseteq \Delta$, bounding an open disc in $\Delta$ including $Z$, such that every edge of $C$ is incident with a region $r$ of $H$ satisfying $\bar{r} \cap Z \ne \emptyset$.}

\textit{Proof.} Let $\mathscr{R}$ be the set of all regions $r$ of $H$ such that $\bar{r} \cap Z \ne \emptyset$. By (5.1) it suffices to show that $\mathscr{R}$ is connected, and this follows easily from the connectedness of $Z$. \fix{We give the details. Since $U(H)$ includes no non-empty open subset of $\Sigma$, every point of $\Sigma$ is in the closure of some region of $H$. In particular $\mathscr{R} \ne \emptyset$, since $Z \ne \emptyset$. To show that $\mathscr{R}$ is connected, choose $r_0 \in \mathscr{R}$, let $\mathscr{R}_1$ be the set of all $r \in \mathscr{R}$ for which there is a sequence from $r_0$ to $r$ as in the definition of connectedness, and let $\mathscr{R}_2 = \mathscr{R} - \mathscr{R}_1$. Suppose that $\mathscr{R}_2 \ne \emptyset$. Let $K_1$, $K_2$ and $K_3$ be the unions of the closures of the members of $\mathscr{R}_1$, of the members of $\mathscr{R}_2$, and of the regions of $H$ not in $\mathscr{R}$, respectively. These sets are closed, since $H$ has only finitely many regions. Every point of $\Sigma$ not in $K_3$ is in $K_1 \cup K_2$, and $Z \cap K_3 = \emptyset$ by the definition of $\mathscr{R}$. We claim that $K_1 \cap K_2 \subseteq V(H)$. For let $x \in \bar{r}_1 \cap \bar{r}_2$, where $r_1 \in \mathscr{R}_1$ and $r_2 \in \mathscr{R}_2$. If $x \notin U(H)$, then the region of $H$ that contains $x$ is an open set meeting $r_1$ and $r_2$, and so $r_1 = r_2$, which is impossible. If $x$ belongs to an edge $e$ of $H$, then $e$ is incident with $r_1$ and $r_2$, and so $r_2 \in \mathscr{R}_1$, which is impossible. This proves our claim.

Next we show that $Z \cap K_1 \cap K_2 = \emptyset$. For suppose that $v \in Z \cap K_1 \cap K_2$. Then $v \in V(H)$ by our claim, and $v \notin K_3$. Since $v \in \Delta - \text{bd}(\Delta)$, $v$ is not in $\text{bd}(\Sigma)$, and so there is an open disc $W \subseteq \Sigma$ containing $v$ and disjoint from $K_3$ and from $V(H) - \{v\}$. Then $W - \{v\}$ is connected. Each point of $W - \{v\}$ is in $K_1 \cup K_2$, since it is not in $K_3$, and it is not in both $K_1$ and $K_2$, by our claim, since it is not in $V(H)$. Since $K_1, K_2$ are closed, it follows that $W - \{v\}$ is disjoint from one of $K_1, K_2$, say from $K_i$. But $v \in K_i$, and so $v \in \bar{r}$ for some region $r$ with $\bar{r} \subseteq K_i$. Then $r$ meets $W$, and $r \cap W \subseteq W - \{v\}$ since $v \in U(H)$. Hence $W - \{v\}$ meets $K_i$, a contradiction. Thus $Z \cap K_1 \cap K_2 = \emptyset$.

Now $Z \subseteq K_1 \cup K_2$, since $Z \cap K_3 = \emptyset$, and $Z$ meets both $K_1$ and $K_2$, since $\mathscr{R}_1$ and $\mathscr{R}_2$ are non-empty subsets of $\mathscr{R}$. Hence $Z$ is the union of the disjoint non-empty sets $Z \cap K_1$ and $Z \cap K_2$, which are closed in $Z$. This is contrary to the connectedness of $Z$, since $Z$ is arc-wise connected. Thus $\mathscr{R}_2 = \emptyset$, and $\mathscr{R}$ is connected.

By (5.1), there is a circuit $C$ of $H$ with $U(C) \subseteq \Delta$, bounding a closed disc $D \subseteq \Delta$ that includes each member of $\mathscr{R}$, such that every edge of $C$ is incident with a member of $\mathscr{R}$. It remains to show that $Z \subseteq D - U(C)$, the open disc bounded by $U(C)$. The set $\Sigma - K_3$ is open, and it includes $Z$. Moreover $\Sigma - K_3 \subseteq D$. For otherwise $(\Sigma - K_3) - D$ is a non-empty open set, and so it contains a point $y \notin U(H)$. Let $r$ be the region of $H$ that contains $y$. Then $r \notin \mathscr{R}$, since $y \notin D$ and every member of $\mathscr{R}$ is a subset of $D$, and so $y \in K_3$, a contradiction. Hence $Z$ is included in the interior of $D$, which is $D - U(C)$. This proves (5.2).}\fixnote{fix:5.2-proof-added}{The paper says only that the connectedness of $\mathscr{R}$ follows easily from that of $Z$, and does not show that the disc given by (5.1) has $Z$ in its interior. The argument above is the one used in the formalization.} $\square$
